% 第5章图：由消元次序(1,4,2,3)得到袋，并入被包含的小袋后得到两袋树分解。
\begin{tikzpicture}[x=1mm,y=1mm,>=Stealth,
  v/.style={circle,draw=FigureStroke,fill=FigureBlueFill,line width=.7pt,minimum size=5.6mm,inner sep=0pt},
  e/.style={FigureStroke,line width=.7pt},
  bag/.style={rectangle,rounded corners=1.8pt,draw=FigureStroke,fill=FigureBlueFill,line width=.7pt,minimum height=6mm,inner sep=1.6pt},
  small/.style={bag,fill=white,draw=FigureMuted},
  sc/.style={font=\figurefont\footnotesize,text=FigureInk},
  note/.style={font=\FigureNoteFont\fontsize{7.5}{9}\selectfont,text=FigureMuted}]
% left: primal graph
\begin{scope}
 \node[v] (a) at (0,8) {$1$}; \node[v] (b) at (9,16) {$2$}; \node[v] (c) at (9,0) {$3$}; \node[v] (d) at (18,8) {$4$};
 \draw[e] (a)--(b) (a)--(c) (b)--(c) (b)--(d) (c)--(d);
 \node[sc] at (9,-10) {(a) 原始图};
\end{scope}
% middle: bags from elimination steps
\begin{scope}[xshift=30mm]
 \node[bag] (B1) at (0,16) {$\{1,2,3\}$};
 \node[bag] (B4) at (0,0) {$\{2,3,4\}$};
 \node[small] (B2) at (20,8) {$\{2,3\}$};
 \node[small] (B3) at (36,8) {$\{3\}$};
 \draw[->,FigureStroke,line width=.7pt] (B1)--(B2) node[midway,above,sc]{$\psi_1$};
 \draw[->,FigureStroke,line width=.7pt] (B4)--(B2) node[midway,below,sc]{$\psi_4$};
 \draw[->,FigureStroke,line width=.7pt] (B2)--(B3) node[midway,above,sc]{$\psi_2'$};
 \node[sc] at (18,-10) {(b) 消元步骤的袋};
\end{scope}
% right: two-bag tree
\begin{scope}[xshift=84mm]
 \node[bag] (L) at (0,16) {$\{1,\textcolor{FigureGreen}{2,3}\}$};
 \node[bag] (R) at (0,0) {$\{\textcolor{FigureGreen}{2,3},4\}$};
 \draw[FigureStroke,line width=.8pt] (L)--(R) node[midway,right,font=\figurefont\footnotesize,text=FigureGreen]{接口$\{2,3\}$};
 \node[sc] at (4,-10) {(c) 并入小袋后};
\end{scope}
\end{tikzpicture}
